% MS4 Final Report — Hyperbolic Embeddings for Hierarchical Style Representation
% Target length: 2000-2500 words of body prose. Title page, headings,
% references, and appendix do NOT count toward the budget. Captions DO.
%
% This file uses the standard article class for portability. To swap in the
% NeurIPS 2026 template later, change \documentclass{article} to
% \documentclass{neurips_2026} (or the version supplied by the course) and
% remove the geometry/title block below; the section structure will carry over.

\documentclass[11pt]{article}

\usepackage[margin=1in]{geometry}
\usepackage{microtype}
\usepackage{amsmath, amssymb, amsthm}
\usepackage{graphicx}
\usepackage{booktabs}
\usepackage{multirow}
\usepackage{caption}
\usepackage{subcaption}
\usepackage{xcolor}
\usepackage[colorlinks=true, linkcolor=blue!50!black, citecolor=blue!50!black, urlcolor=blue!50!black]{hyperref}
\usepackage[numbers,sort&compress]{natbib}
\usepackage{authblk}
\usepackage{lipsum} % remove before submission

\title{Hyperbolic Embeddings for Hierarchical Style Representation: \\
  An Empirical Study on WikiArt}

\author[1]{Peter Flo}
\author[1]{Luca Grossmann}
\author[1]{Valerie Wang}
\affil[1]{Harvard University, COMPSCI 2090B (Spring 2026)}
\date{\textbf{Canvas Project \#150} \\ \today}

\renewcommand{\abstractname}{Abstract}

\begin{document}

\maketitle

\begin{abstract}
% ~150 words. Final pass after results are settled.
We investigate whether Poincar\'e-ball embeddings recover the hierarchical
structure of artistic styles better than matched Euclidean embeddings. Using
frozen CLIP ViT-B/16 features over 81{,}000 WikiArt paintings spanning 27
styles, we train geometry-paired prototype classifiers across embedding
dimensions $d \in \{2,\dots,64\}$, curvatures $c \in \{0.1, 0.3, 1, 3\}$, and
seeds, both with and without a hierarchy-aware regularizer that pushes the
prototype distance matrix toward a hand-built art-history tree. We find
\textbf{[headline result placeholder]}. Tree-variant ablations with
chronological and flat ground-truth trees confirm that the conclusion is
\textbf{[robust / fragile]} to the specific choice of reference hierarchy.
Code, sweep CSVs, and figures are available at the project repository.
\end{abstract}

% =================================================================
\section{Background and Motivation} % \approx 300 words. Owner: Luca.
% =================================================================
% Why hierarchical structure in style? Why hyperbolic? Connect to broader
% representation-learning questions; mention CLIP frozen features as a
% practical foundation; flag the existing hyperbolic-embeddings literature
% (Nickel \& Kiela 2017, Ganea et al. 2018, Khrulkov et al. 2020).

\lipsum[1]

% =================================================================
\section{Problem Statement} % \approx 100 words. Owner: shared.
% =================================================================
% Refined research question, target/output, unit of observation.

We ask: \emph{do Poincar\'e-ball embeddings trained on frozen CLIP features
recover the WikiArt style hierarchy more faithfully than Euclidean embeddings
of matched dimensionality and training budget?} The unit of observation is
one painting; the target during training is its 27-way style label, with the
hand-built lineage tree available either as an evaluation reference or as an
auxiliary loss. We report classification accuracy, prototype-tree distance
correlation, hierarchical clustering agreement, and local retrieval at the
sibling and cousin level.

% =================================================================
\section{Data and Exploratory Analysis} % \approx 250 words. Owner: Valerie.
% =================================================================
% Pull only the EDA findings that drove modeling decisions:
% - Class imbalance (drives balanced-acc reporting + prototype classifier).
% - Hierarchy non-uniqueness (drives the tree-variant ablation in Sec. 6.3).
% - Image preprocessing pipeline summary (CLIP-style normalization, 224px).

\lipsum[2]

% =================================================================
\section{Methods} % \approx 600 words. Owner: Peter.
% =================================================================

\paragraph{Architecture.} A two-layer MLP with GELU and dropout maps frozen
CLIP ViT-B/16 features ($\mathbb{R}^{512}$) to a $d$-dimensional embedding.
A geometry-specific head produces the final embedding: identity for the
Euclidean case, exponential map at the origin for the Poincar\'e ball,
\(\exp_0^c(u) = \tanh(\sqrt{c}\,\|u\|) \, u / (\sqrt{c}\,\|u\|)\). A
prototype classifier with one learnable point per style yields logits as
$-d(z, p_k)$, where $d(\cdot,\cdot)$ is Euclidean or hyperbolic accordingly.

\paragraph{Geometry, distance, and optimization.} The Poincar\'e distance
with curvature $c > 0$ is
\[
  d^c(x, y) = \frac{1}{\sqrt{c}}\,\operatorname{arcosh}\!\left(
    1 + \frac{2c\,\|x-y\|^2}{(1-c\|x\|^2)(1-c\|y\|^2)}
  \right).
\]
Linear layers do not preserve the ball, so we cannot use a plain logit head;
the prototype classifier is a natural alternative because logits are
distances. The MLP head is optimized with Adam; the hyperbolic prototypes
are optimized with Riemannian Adam from \texttt{geoopt}, which projects each
update back onto the manifold.

\paragraph{Hierarchy-aware loss.} The MS3 baseline trains with cross-entropy
only and uses the hand-built tree only at evaluation time. For MS4 we add a
regularizer that pushes the prototype-prototype distance matrix toward the
tree distance matrix:
\[
  \mathcal{L} = \mathcal{L}_{CE} + \lambda \cdot \frac{1}{|\mathcal{P}|}
    \sum_{(i,j) \in \mathcal{P}} \!\left(
      \frac{d(p_i, p_j)}{\bar{d}} - \frac{T_{ij}}{\bar{T}}
    \right)^{\!\!2}.
\]
The mean-normalization makes the term scale-invariant, so the geometry
chooses its own absolute scale rather than fighting the classification
objective for embedding magnitude. $\lambda = 0$ recovers MS3.

\paragraph{Ground-truth trees.} The default tree is a hand-built lineage
hierarchy following standard art-history references (Renaissance branches
into Baroque, etc.; see Appendix~\ref{app:tree}). We additionally evaluate
two variants: a \emph{chronological} two-level era grouping
(Renaissance / Baroque-Rococo / 19th century / Early 20th / Modern /
Non-Western) and a \emph{flat} tree (every leaf is a direct child of the
root, used as a null).

\paragraph{Training details.} All runs use the supplied 70/30 train/val
split ($\approx$57k / 24k images), batch size 4096, Adam $\eta=10^{-3}$ for
the head and $10^{-2}$ for the prototypes, weight decay $10^{-4}$, dropout
0.1, gradient clipping at norm 1.0, 30 epochs (training is fast because the
CLIP features are pre-cached). Each headline configuration is run with three
seeds. The full grid (Phases 1–3, see Section~\ref{sec:results}) totals
$\approx 200$ runs.

\paragraph{Evaluation suite.} We report top-1, top-5, and balanced
accuracy; prototype-tree and class-center-tree Spearman correlation;
multiplicative tree distortion (mean and worst-case); dendrogram-cluster
F1; sibling and cousin recall@$k$ for $k \in \{5, 10\}$; and Fr\'echet-mean
nearest-prototype accuracy on Cubism as an interpolation-stability probe.
Implementation details are in \texttt{scripts/eval.py}.

% =================================================================
\section{Results and Interpretation}\label{sec:results} % \approx 700 words. Owner: Peter / Valerie.
% =================================================================

\subsection{Phase 1 — Scaling sweep}
% TODO: insert figure: Phase 1 scaling plot (top-1, tree spearman, sibling@5
% vs d, by geometry; hyperbolic uses best curvature per d).
\begin{figure}[t]
  \centering
  \includegraphics[width=\linewidth]{figures/phase1_scaling.pdf}
  \caption{Phase 1 — metrics versus embedding dimension $d$. Hyperbolic
  bars show the best curvature per $d$; error bars span seed standard
  deviation. \textbf{[Insert narrative: does hyperbolic close the gap with
  $d$? At what $d$? On which metric?]}}
  \label{fig:phase1}
\end{figure}

\subsection{Phase 2 — Hierarchy-aware loss}
% TODO: insert figure: metrics vs lambda for both geometries at d=8.
\begin{figure}[t]
  \centering
  \includegraphics[width=\linewidth]{figures/phase2_lambda.pdf}
  \caption{Phase 2 — top-1, class-center / tree Spearman, tree distortion,
  and sibling recall@5 versus the hierarchy-aware loss weight $\lambda$ at
  $d = 8$, $c = 1$. \textbf{[Insert narrative: does hyperbolic catch up on
  global metrics once trees are used at training time?]}}
  \label{fig:phase2}
\end{figure}

\subsection{Phase 3 — Tree-variant ablation}
% TODO: insert table comparing default / chronological / flat training trees.

\subsection{Headline comparison}
% TODO: insert summary table — best Euclidean vs best Hyperbolic across
% phases, with seed standard deviations.

\subsection{Qualitative analysis}
% TODO: Poincare-disk plot at d=2 (prototypes + sample embeddings),
% nearest-neighbor grid, confusion matrix block-ordered by tree.

\lipsum[3]

% =================================================================
\section{Conclusions and Discussion} % \approx 250 words.
% =================================================================
% What we can claim, what we can't, what we'd test next.

\lipsum[4]

% =================================================================
\section{Broader Impact} % \approx 100 words.
% =================================================================
% Cultural-heritage applications, risks of imposing a specific art-historical
% tree, dataset biases (Western-canon-heavy), de-skilling concerns.

\lipsum[5]

\bibliographystyle{plainnat}
\bibliography{references}

\appendix

\section{Style hierarchy used as the default tree}\label{app:tree}
% Drop the dict from scripts/hierarchy.py here as a verbatim listing.

\section{Full sweep results}
% Optional: full results CSV summary tables.

\end{document}
